\section{背景与动机}
实验谱图是分子结构解析中最重要的证据来源之一：IR/Raman 反映键振动与官能团环境，NMR 提供局部化学环境与连接关系线索，MS 则对分子式、碎片与分子量范围给出约束。近年的数据集与模型工作表明，多模态联合输入能够显著增强结构解析上限，而前向谱图预测模型已经具备为候选分子提供快速物理一致性评分的能力 \citep{Alberts2024MultimodalDataset,RocabertOriols2025VibraCLIP,Zou2023DetaNet}。

但将谱图直接翻译为 SMILES 仍存在三类关键局限。第一，SMILES 是人为序列，不是分子本体表示；同一分子可对应多种等价书写方式。第二，谱图到结构本质上是一谱多解的逆问题，单一序列输出会压缩不确定性表达。第三，谱图的物理来源依赖分子图、三维几何与张量响应；若忽略几何与前向验证，模型容易在实验域噪声下失去可解释性。以 IR/Raman 到结构的双向翻译工作为例，直接“谱图到结构”的范式虽有进展，但实验域性能、结构歧义与候选可验证性依然是主要瓶颈 \citep{Hu2025TranSpec,Alberts2024IRStructure}。

因此，本项目不把 demo 目标设定为“训练一个完整的大模型”，而是构造一条更符合科学反演逻辑的工作流：先从多模态谱图中提取结构证据，再生成或检索候选分子，最后通过前向谱图一致性与化学有效性完成重排。这样的设计更适合 DP Technology / Deep Potential 的 AI for Science 场景：它强调可追踪、可验证、可扩展，而不是只追求端到端生成结果的表面准确率。

\begin{keybox}[为什么适合作为实习 Demo]
\begin{itemize}
  \item 问题清晰：谱图到结构反演是典型的科学逆问题；
  \item 路径可落地：可先做最小闭环，再逐步引入更强生成与几何模块；
  \item 结果可验证：Top-K 输出、前向谱图一致性与 trace log 均可直接检查；
  \item 技术外延强：可自然扩展至图扩散、几何深度学习与物理先验融合。
\end{itemize}
\end{keybox}
